% Code areas by number of compounds, stacked by the digits of their name's range (registry/areas.csv.gz).
% Data: area-sizes.dat, written by tools/paper-figures.py.
\documentclass[tikz,border=3pt]{standalone}
\usepackage{fontspec}
\setmainfont{LibertinusSans}[Extension=.otf,UprightFont=*-Regular,BoldFont=*-Bold,ItalicFont=*-Italic]
\usepackage{unicode-math}
\setmathfont{LibertinusMath-Regular.otf}
\usepackage{pgfplots}
\pgfplotsset{compat=1.18}
% English by default; the French paper compiles <name>-fr.tex, which sets \fr to 1 and inputs this file.
\providecommand{\fr}{0}
\newcommand{\T}[2]{\ifnum\fr=1 #2\else #1\fi}
\ifnum\fr=1 \pgfkeys{/pgf/number format/.cd, use comma, 1000 sep={\,}}\fi
\definecolor{ink}{HTML}{1F5FA8}\definecolor{flaggreen}{HTML}{009E60}\definecolor{yellow}{HTML}{FCD116}
\definecolor{flagblue}{HTML}{3A75C4}\definecolor{amber}{HTML}{D9A400}
\begin{document}
\begin{tikzpicture}
\begin{axis}[ybar stacked, bar width=17pt, width=12cm, height=6.2cm, ymin=0, ymax=1000,
  xtick=data, xticklabels={1, 2–3, 4–9, 10–31, 32–99, 100–316, 317–999, \T{1k–3.2k}{1–3,2 k}, \T{3.2k–10k}{3,2–10 k}, \T{≥10k}{≥ 10 k}},
  x tick label style={font=\scriptsize}, y tick label style={font=\footnotesize},
  xlabel={\T{Compounds in the area}{Concessions dans la zone}}, ylabel={\T{Areas}{Zones}}, label style={font=\small},
  legend style={at={(0.98,0.97)}, anchor=north east, font=\footnotesize, draw=gray!50},
  legend cell align=left, axis line style={gray}, ymajorgrids, grid style={gray!20},
  enlarge x limits=0.06]
\addplot[fill=flaggreen, draw=none] table[x=bin, y=d3] {area-sizes.dat};
\addplot[fill=flagblue, draw=none] table[x=bin, y=d4] {area-sizes.dat};
\addplot[fill=amber, draw=none] table[x=bin, y=d5] {area-sizes.dat};
\legend{\T{name with 3 digits}{nom à 3 chiffres}, \T{4 digits}{4 chiffres}, \T{5 digits}{5 chiffres}}
\end{axis}
\end{tikzpicture}
\end{document}
